\src{187}
\result{COROLLARY}{} One has $H^{i}(X,F(-m))=0$ for $m$ large if and only if
$E^{n-i}(F)=0$.

Recall that the $E^{j}(F)$ do not in fact depend on the projective embedding
under consideration. The condition in the corollary is purely local.
Consequently, if it holds for $F$, it holds for every sheaf locally isomorphic
to a sheaf $F^{n}$. In particular, if this condition holds for $\Osh_X$, it
holds for every locally free coherent algebraic sheaf. This is the case, for
example, for every $i<n$ if $X$ is nonsingular; for $i=0$ if no component of
$X$ is reduced to a single point; and for $i=0,1$ if $X$ is normal and all its
components have dimension $>1$ (cf. [3]).\footnote{The printed text says here
``$S$ is normal''; the paragraph involves only the variety $X$.} For this to
hold for every $i<k$, it is necessary and sufficient, by definition, that the
local rings $\Osh_x$ ($x\in X$) have ``homological codimension'' $\geq k$ (see
[4] for this notion). If $k=n$, Theorem 3 bis says that this means the duality
theorem is true for $X$, i.e. that (6.5) is an isomorphism for every $p$ and
every $F$. Many equivalent conditions on the local rings $\Osh_x$ can be given
for this to hold (Nagata), for example that they satisfy the Cohen--Macaulay
equidimensionality theorem. It holds, for example, if $X$ is locally a
``complete intersection'' in a nonsingular ambient variety.

\fsection{7}{Poincaré duality}

Let $X$ be a nonsingular projective variety of dimension $n$. Then
\[
  H^{**}(X)=H^{*}(X,\Omega_X^{*})
\]
is a finite-dimensional anticommutative bigraded algebra, which we grade by
total degree, so that $H^{p,q}(X)=H^{p}(X,\Omega_X^{q})$ has degree $p+q$. The
degrees of $H^{*}(X)$ range from $0$ to $2n$. By Theorem 2 and the corollary to
Theorem 3, $H^{*}(X)$ is a ``Poincaré algebra'' of dimension $2n$: that is,
$H^{2n}(X)$ is equipped with an isomorphism to the base field $k$, and the
product
\[
  H^{m}(X)\times H^{2n-m}(X)\longrightarrow H^{2n}(X)=k
\]
is a duality between $H^{m}(X)$ and $H^{2n-m}(X)$. Moreover, if $Y$ is another
nonsingular projective variety, the Künneth formula for coherent algebraic
sheaves gives
\begin{equation}\tag{7.1}
  H^{*}(X\times Y)=H^{*}(X)\otimes H^{*}(Y),
\end{equation}
an isomorphism compatible with the Poincaré-algebra structures. Furthermore,
as a commutative algebra, $H^{*}(X)$ is a contravariant functor of $X$: a
morphism $f:Y\longrightarrow X$ defines in the evident way a homomorphism of
graded algebras

